\section{Preliminaries}\label{prel.sec}

Let $[n]\eqdef \{1,\ldots,n\}$, and let
$\calD = (D_1,\ldots,D_m)$ be a list of $m$ distributions,
where $D_i : [n] \to [0,1]$ and $\sum_{j=1}^n D_i(j) =1$
for every $1 \leq i \leq m$.
For a vector $\vbf = (v_1,\ldots,v_n) \in \mathbb{R}^n$,
let $\|\vbf\|_1 = \sum_{i=1}^n |v_i|$ denote the
$\ell_1$ norm of the vector $v$.

%
For a property  $\calP$ of lists of distributions and
$0 \leq \eps \leq 1$, we say that
$\calD$ is {\em $\eps$-far from\/} (\emph{having}) $\calP$ if
$({1}/{m}) \sum_{i=1}^m \|D_i - D^*_i\|_1 > \eps$
for every list $\calD^* = (D^*_1,\ldots,D^*_m)$
that has the property $\calP$. Note that
the statistical distance (also known as the total variation distance)
between two discrete distributions, $D$ and $D^*$ over $[n]$,
that is,
$\max_{S\subseteq [n]} \left| \sum_{i\in S} D(i) -\sum_{i\in S} D^*(i)\right|$,
 equals $\|D - D^*\|_1/2$.

Given a distance parameter $\eps$,
a testing algorithm for a property $\calP$
%
should distinguish between the case that $\calD$ has the property
$\calP$ and the case that it is $\eps$-far from $\calP$. We consider
two models within which this task is performed.
\begin{description}
\item[The Query Model]
In this model the testing algorithm may indicate an index $1 \leq i \leq m$ of
its choice and it gets an independent sample $j$ distributed according to $D_i$.
\item[The Sampling Model]
In this model the algorithm cannot select (``query'') a distribution of
its choice. Rather, it may obtain a pair $(i,j)$ where $i$ is selected
uniformly (we refer to this as the \emph{Uniform} sampling model)
and $j$ is an independent sample distributed according to $D_i$.

We also consider a generalization in which there is an underlying
weight vector $\wbf = (w_1,\ldots,w_m)$ (where $\sum_{i=1}^m w_i = 1$),
and the distribution $D_i$ is selected according to $\wbf$.
In this case the notion of ``$\eps$-far'' %
needs to be modified accordingly.
%
We shall say  %
that $\calD$ is $\eps$-far from $\calP$ \emph{with respect to $\wbf$}
if $\sum_{i=1}^m w_i\cdot \|D_i - D^*_i\|_1 > \eps $
for every list $\calD^* = (D^*_1,\ldots,D^*_m)$
that has the property $\calP$.

We consider two variants of this non-uniform model:
The \emph{known-weights} sampling model, in which $\wbf$ is known
to the algorithm, and the \emph{unknown-weights} sampling model
in which $\wbf$ is known.
%
%
%
%
%
\end{description}

%
%
%
%
%
%

A main focus of this work is on the following property.
We shall say that a list $\calD = (D_1 \ldots D_m)$ of $m$ distributions over $[n]$
belongs to $\calPidmn$ (or has the property $\calPidmn$) if $D_i = D_{i'}$ for all
$1 \leq i, i' \leq m$.

\section{A lower bound for testing equivalence 
in the uniform sampling model}\label{id-lb.sec}
%
In this section we prove the following theorem:
\begin{theorem}
\label{lb1.thm}
Any testing algorithm for the property  $\calPidmn$  in the uniform sampling model
for every $\eps \leq 1/20$ and for
 $n > c m\log m$ where $c$ is some sufficiently large
constant,
%
%
requires $\Omega(n^{2/3}m^{1/3})$ samples.
%
%
%
%
\end{theorem}
The proof of \expref{Theorem}{lb1.thm} consists of two parts. The first
is a general theorem (\expref{Theorem}{wish.thm})
concerning testing  symmetric properties of lists of distributions.
%
%
This theorem extends
 a lemma of Valiant~\cite[Lem. 4.5.4]{valphd}
(which leads to what Valiant refers to as
the ``Wishful Thinking Theorem''). The second part is a construction of
two lists of distributions to which \expref{Theorem}{wish.thm} is applied.
Our analysis uses a technique called \emph{Poissonization}~\cite{szpankowski}
(which was used in the past in the context of lower bounds for testing and estimating
properties of distributions in~\cite{RRSS,val,valphd}),
and hence we first introduce some preliminaries concerning Poisson distributions.
%
We later provide some intuition regarding the benefits of Poissonization.




\subsection{Preliminaries concerning Poisson distributions}
\label{prel-poiss.subsec}
\sloppy
For a positive real number $\lambda$,
the Poisson distribution $\poi(\lambda)$ takes the value
$x \in \mathbb{N}$ (where $\mathbb{N} = \{0,1,2,\ldots\}$) with
probability $\poi(x;\lambda) = \eee^{-\lambda} \lambda^x/x!$.
The expectation and variance of $\poi(\lambda)$ are both $\lambda$.
For $\lambda_1$ and $\lambda_2$ we shall use the following bound
on the $\ell_1$ distance between the corresponding Poisson distributions
(for a proof see for example~\cite[Claim A.2]{RRSS}):
\BEQ
\|\poi(\lambda_1) - \poi(\lambda_2)\|_1 \leq 2|\lambda_1 - \lambda_2|\,.
\label{poi-stat.eq}
\EEQ

For a vector $\vec{\lambda} = (\lambda_1,\ldots,\lambda_d)$ of positive
real numbers, the corresponding \emph{multivariate} Poisson distribution
$\poi(\vec{\lambda})$ is the product distribution
$\poi(\lambda_1)\times\dots\times \poi(\lambda_d)$. That is,
$\poi(\vec{\lambda})$ assigns each vector
$\vec{x} = x_1\ldots,x_d \in \mathbb{N}^d$ the probability
$\prod_{i=1}^d \poi(x_i;\lambda_i)$.

We shall sometimes consider vectors $\vec{\lambda}$ whose coordinates are indexed
by vectors  $\vec{a} = (a_1,\ldots,a_m) \in \mathbb{N}^m$, and will use
$\vec{\lambda}(\vec{a})$ to denote the coordinate of $\vec{\lambda}$
that corresponds to $\vec{a}$. Thus, $\poi(\vec{\lambda}(\vec{a}))$
is a univariate Poisson distribution.
With a slight abuse of notation, for a subset $I \subseteq [d]$
(or $I \subseteq \mathbb{N}^m$),
we let $\poi(\vec{\lambda}(I))$ denote the multivariate Poisson
distributions restricted to the coordinates of $\vec{\lambda}$ in $I$.

For any two $d$-dimensional vectors $\vec{\lambda}^+ = (\lambda^+_1,\ldots,\lambda^+_d)$ and
$\vec{\lambda}^- = (\lambda^-_1,\ldots,\lambda^-_d)$ of positive real values, we get from the proof of~\cite[Lemma 4.5.3]{valphd} that
%

\BEQN
\left\|\poi(\vec{\lambda}^+) - \poi(\vec{\lambda}^-) \right\|_1
   \leq
    \sum_{j=1}^d \left\|\poi(\lambda^+_j) - \poi(\lambda^-_j) \right\|_1 \,. \label{lam.eq}
\EEQN







%
For our purposes we generalize equation~(\ref{lam.eq}) in the following lemma.

\begin{lemma} \label{disval.lem}
For any two $d$-dimensional vectors $\vec{\lambda}^+ = (\lambda^+_1,\ldots,\lambda^+_d)$ and
$\vec{\lambda}^- = (\lambda^-_1,\ldots,\lambda^-_d)$ of positive real values, and for
any partition $\{I_i\}_{i=1}^{\ell}$ of $[d]$,
%
%
%

\BEQN
\left\|\poi(\vec{\lambda}^+) - \poi(\vec{\lambda}^-) \right\|_1
   \leq
    \sum_{i=1}^\ell \left\|\poi(\vec{\lambda}^+(I_i))
      - \poi(\vec{\lambda}^-(I_i)) \right\|_1 \,.
\EEQN








%
\end{lemma}

%
\def\DisvalLemPrf{
Let $\{I_i\}_{i=1}^{\ell}$ be a partition of $[d]$, let $\vec{i}$ denote $(i_1,\ldots i_d)$, by the triangle inequality we have that for every $k\in [\ell]$,
%

\begin{eqnarray}
\left|\poi(\vec{i}\,;\,\vec{\lambda}^+) - \poi(\vec{i}\,;\,\vec{\lambda}^-)\right|
   &=& \Big|\prod_{j\in[d]} \poi(i_j;\lambda_j^+) - \prod_{j\in[d]} \poi(i_j;\lambda_j^-)\Big| \\
   &\leq& \Big|\prod_{j\in[d]} \poi(i_j;\lambda_j^+) - \prod_{j\in[d]\setminus I_k} \poi(i_j;\lambda_j^+) \prod_{j\in I_k} \poi(i_j;\lambda_j^-)\Big| \\ &&\;+\;
    \Big|\prod_{j\in[d]\setminus I_k} \poi(i_j;\lambda_j^+) \prod_{j\in I_k} \poi(i_j;\lambda_j^-) - \prod_{j\in[d]} \poi(i_j;\lambda_j^-)\Big|\,.
\end{eqnarray}










%
Hence, we obtain that
%

\begin{eqnarray}
\left\|\poi(\vec{\lambda}^+) - \poi(\vec{\lambda}^-) \right\|_1
   &=& \sum_{\vec{i} \in \mathbb{N}^d} \left|\poi(\vec{i}\,;\,\vec{\lambda}^+) - \poi(\vec{i}\,;\,\vec{\lambda}^-)\right| \\
   &\leq& \left\|\poi(\vec{\lambda}^+(I_k))
      - \poi(\vec{\lambda}^-(I_k)) \right\|_1 \\
   && + \left\|\poi(\vec{\lambda}^+([d]\setminus I_k))
      - \poi(\vec{\lambda}^-([d] \setminus I_k)) \right\|_1 \; .
\end{eqnarray}












%
Thus, the lemma follows by induction on $\ell$.
}
%

%

\BPF
\DisvalLemPrf
\EPF

%

%
\noindent
We shall also make use of the following lemma.
\begin{lemma} \label{dis.lem}
%
For any two $d$-dimensional vectors $\vec{\lambda}^+ = (\lambda^+_1,\ldots,\lambda^+_d)$ and
$\vec{\lambda}^- = (\lambda^-_1,\ldots,\lambda^-_d)$ of positive real values,
%

\BEQN
\left\|\poi(\vec{\lambda}^+) - \poi(\vec{\lambda}^-) \right\|_1
   \leq 2\sqrt{2\sum_{j=1}^d \frac{(\lambda_j^- - \lambda_j^+)^2}{\lambda_j^-}}\,.
\EEQN






%
\end{lemma}


%
\def\DisLemPrf{
In order to prove the lemma we shall use the \emph{KL-divergence} between
distributions. 
%
For
two distributions $p_1$ and $p_2$ over a domain $X$,
we set %
\[
D_{\rm KL}(p_1 \|\, p_2) \eqdef \sum_{x \in X} p_1(x)\cdot
\ln\frac{p_1(x)}{p_2(x)}\,.
\]
Let $\vec{\lambda}^+ = (\lambda_1^+ \ldots,\lambda_d^+)$,
 $\vec{\lambda}^- = (\lambda_1^- \ldots,\lambda_d^-)$ and let $\vec{i}$ denote $(i_1,\ldots i_d)$.
We have that
%

\begin{eqnarray}
\ln \frac{\poi(\vec{i}\,;\,\vec{\lambda}^+)}
                  { \poi(\vec{i}\,;\,\vec{\lambda}^-)}
&=& \sum_{j=1}^d \ln\left(\eee^{\lambda_j^- -\lambda_j^+}
            \left(\lambda_j^+/\lambda_j^-\right)^{i_j}\right) \\
&=& \sum_{j=1}^d \Big((\lambda_j^- - \lambda_j^+) + i_j \cdot \ln(\lambda_j^+/\lambda_j^-)\Big) \\
&\leq& \sum_{j=1}^d \Big((\lambda_j^- - \lambda_j^+) + i_j \cdot
         (\lambda_j^+/\lambda_j^- - 1)\Big) \;,
\end{eqnarray}











%
where in the last inequality we used the fact that
$\ln x \leq x -1$ for every $x>0$.
Therefore, we obtain that
\allowdisplaybreaks
%

\BEQN
D_{\rm KL}\left(\poi(\vec{\lambda}^+)\|\,\poi(\vec{\lambda}^-)\right)
 &=&
    \sum_{\vec{i} \in \mathbb{N}^d} \poi(\vec{i}\,;\,\vec{\lambda}^+) \cdot
         \ln \frac{     \poi(\vec{i}\,;\,\vec{\lambda}^+)}
                  { \poi(\vec{i}\,;\,\vec{\lambda}^-)}\\
&\leq& \sum_{j=1}^d \Big((\lambda_j^- - \lambda_j^+) + \lambda_j^+ \cdot
         (\lambda_j^+ /\lambda_j^- - 1)\Big) \label{eq:kl:1}\\
&=& \sum_{j=1}^d \frac{(\lambda_j^- - \lambda_j^+)^2}{\lambda_j^-} \;,
\EEQN
where in equation~(\ref{eq:kl:1}) we used the facts that $\sum_{i\in \mathbb{N}} \poi(i;\lambda) = 1$ and
$\sum_{i\in \mathbb{N}} \poi(i;\lambda)\cdot i = \lambda$.
















%
The $\ell_1$ distance is related to the KL-divergence  by
\[
\|D-D'\|_1 \leq 2\sqrt{2D_{\rm KL}\left(D\|\,D'\right)}
\]
%
and thus we obtain the lemma.
}
%
%

\BPF
\DisLemPrf
\EPF

%
%
\def\PoiExpLem{

\noindent
%




\medskip\noindent
The next lemma bounds the probability that a Poisson random variable is significantly
smaller than its expected value.

\begin{lemma} \label{l:poiexp}
Let $X \sim \poi(\lambda)$, then,
\BEQ
\Pr[X < \lambda/2] < (3/4)^{\lambda/4}\,.
\EEQ
\end{lemma}
\BPF
Consider the matching between $j$ and $j + \lambda/2$ for every
$j = 0,\ldots,\lambda/2 - 1$. We consider the ratio between $\poi(j;\lambda)$
and $\poi(j+\lambda/2;\lambda)$:
%

\begin{eqnarray}
\frac{\poi(j+\lambda/2;\lambda)}{\poi(j;\lambda)}
&=&
  \frac{\eee^{-\lambda}\cdot \lambda^{j+\lambda/2} / (j+\lambda/2)!}
       {\eee^{-\lambda}\cdot \lambda^{j}/j!} \\
&=& \frac{\lambda^{\lambda/2}}{(j+\lambda/2)(j+\lambda/2 -1)\cdots(j+1)} \\
&=& \frac{\lambda}{j+\lambda/2}\cdot \frac{\lambda}{j+\lambda/2 -1}\cdots \frac{\lambda}{j+1}\\
&\geq& \frac{\lambda}{\lambda -1}\cdot \frac{\lambda}{\lambda -2}
               \cdots \frac{\lambda}{\lambda/2} \\
&>& \left(\frac{\lambda}{(3/4)\lambda}\right)^{\lambda/4} \\
&=& (4/3)^{\lambda/4}\,.
\end{eqnarray}














%
This implies that
%

\BEQN
\Pr[X < \lambda/2] &=& \frac{\Pr[X < \lambda/2]}{\Pr[\lambda/2 \leq X < \lambda]}
                             \cdot \Pr[\lambda/2 \leq X < \lambda] \\
&<& \frac{\Pr[X < \lambda/2]}{\Pr[\lambda/2 \leq X < \lambda]} \\
&<& (3/4)^{\lambda/4} \;,
\EEQN









%
and the proof is completed.
\EPF
}
%

%

\PoiExpLem

%

The next two notations will play an important technical role in our analysis.
For a list of distributions $\calD = (D_1 \ldots D_m)$, an integer $\w$ and
a vector $\vec{a} = (a_1,\ldots,a_m) \in \mathbb{N}^m$, let
\BEQ
p^{\calD,\w}(j;\vec{a}) \eqdef \prod_{i=1}^m \poi(a_i;\w\cdot D_i(j))\,.
\label{calDi-def.eq}
\EEQ
That is, for a fixed choice of a domain element $j\in[n]$,
consider performing $m$ independent trials, one for each distribution $D_i$, where in trial $i$
we select a non-negative integer according to the Poisson distribution
$\poi(\lambda)$ for $\lambda = \w\cdot D_i(j)$.
Then $p^{\calD,\w}(j;\vec{a})$ is the probability of the joint event
that we get an outcome of $a_i$ in trial $i$, for each $i \in [m]$.
Let $\vec{\lambda}^{\calD,\w}$ be a vector whose coordinates are indexed by
all $\vec{a} \in \mathbb{N}^m$, such that
\BEQ
\vec{\lambda}^{\calD,\w}(\vec{a})=\sum_{j=1}^n p^{\calD,\w}(j;\vec{a})\,.
\label{vecl-calD-def.eq}
\EEQ
That is, $\vec{\lambda}^{\calD,\w}(\vec{a})$ is the expected number of
times we get the joint outcome $(a_1,\ldots,a_m)$ if we perform the
probabilistic process defined above independently for every
$j\in[n]$. In the next subsection we show that the quantity $\vec{\lambda}^{\calD,\w}(\vec{a})$ characterizes the view of the tester,
provided that it tests a symmetric property, and thus a lower bound
can be proved via a construction of two instances of collections that
agree on this quantity while one instance has the property and
the other is far from having the property.

\subsection{Testability of symmetric properties of lists of distributions}
\label{wish.subsec}

%

In this subsection we prove the following theorem (which is used to prove
\expref{Theorem}{lb1.thm}).




%

\begin{theorem} \label{wish.thm}
%
%
Let $\calD^+$ and $\calD^-$
be two lists of $m$ distributions over $[n]$, all of whose frequencies
are at most $\frac{\delta}{\w\cdot m}$ where $\w$ is some positive integer and
%
%
%
%
$0< \delta \leq 1/2$.  If %
\BEQ
\left\|\poi\left(\vec{\lambda}^{\calD^+,\w}\right) -
    \poi\left(\vec{\lambda}^{\calD^-,\w}\right) \right\|_1
   < \frac{31}{60} - \frac{352\delta}{5}\;,
\label{poi-poi.eq}
\EEQ
 then testing in the uniform sampling model
 any symmetric property of distributions  such that
 $\calD^+$ has the property, while
$\calD^-$ is $\Omega(1)$-far  from having the property
requires $\Omega(\w\cdot m)$ samples.
%
\end{theorem}
%
%
\paragraph{A high-level discussion of the proof of \expref{Theorem}{wish.thm}}
For an element $j \in [n]$ and a distribution $D_i$,
$i \in[m]$, let $\alpha_{i,j}$ be the number of times the pair $(i,j)$ appears in the sample
(when the sample is selected according to some sampling model).
Thus $(\alpha_{1, j}, \ldots, \alpha_{m, j})$ is the \emph{sample histogram} of the element $j$.
The histogram of the elements' histograms is called the
\emph{fingerprint} of the sample. That is, the fingerprint indicates,
for every $\vec{a} \in \mathbb{N}^m$, the number of elements $j$ such
that $(\alpha_{1, j}, \ldots, \alpha_{m, j}) = \vec{a}$.
\sloppy
As shown in~\cite{BFRSW-long},
%
 when testing symmetric properties of distributions, it can be assumed without loss of
generality that the testing algorithm is provided only with the
fingerprint of the sample.
%
%
Furthermore, since the number, $n$, of elements is fixed,
it suffices to give the tester the fingerprint
of the sample without the $\vec{0} = (0,\ldots,0)$ entry.

For example, consider the distributions $D_1$ and $D_2$ over $\{1,2,3\}$ such that $D_1[j] = 1/3$ for every $j\in \{1,2,3\}$, $D_2[1]=D_2[2] = 1/2$ and $D_2[3] = 0$. Assume that we sample $(D_1, D_2)$ four times, according to the uniform sampling model and we get the samples $(1,1), (2,1), (2,2), (1,3)$, where the first coordinate
denotes the distribution, and the second coordinate denotes the element.
Then the sample histogram of element $1$ is $(1,1)$ because $1$ was selected once by $D_1$ and once by $D_2$. For the elements $2,3$ we have the sample histograms $(0,1)$  and $(1,0)$, respectively. The fingerprint of the sample is $(0,1,1,0,1,0,0,\ldots)$ for the following order of histograms: $((0,0), (0,1), (1,0), (2,0) (1,1), (0,2), (3,0), \ldots)$.

%
\smallskip
In order to prove \expref{Theorem}{wish.thm}, we would like to show that
the distributions of the fingerprints when the sample is generated according
to $\calD^+$ and when it is generated according to $\calD^-$ are similar,
for a sample size that is below the lower bound stated in the theorem.
%
%
For each choice of
element $j\in [n]$ and a distribution $D_i$, %
the number of times  the sample $(i,j)$ appears, \ie, $\alpha_{i,j}$,
depends on the number of times the other samples appear simply because
the total number of samples is fixed.
Furthermore,
%
for each histogram $\vec{a}$,
the number of elements with sample histogram identical to $\vec{a}$
is dependent on the number of times the other histograms appear,
because the number of samples is fixed.
For instance, in the example above, if we know that we have the
histogram $(0,1)$ once and the histogram $(1,1)$ once,
 then we know that third histogram cannot be $(2,0)$. In addition,
it is dependent because the number of elements is fixed.

%
%
We thus see that the distribution of the fingerprints is rather difficult to
analyze (and therefore it is difficult
 to bound the statistical distance between two different such distributions).
Therefore, we would like to break as much of the above dependencies.
To this end we %
define a slightly different process for generating
the samples that involves \emph{Poissonization}~\cite{szpankowski}.
In the Poissonized process the number of samples we take from each distribution $D_i$,
denoted by $\w'_i$, is distributed according to the Poisson distribution.
We prove that, while the overall number of samples the Poissonized process
takes is bigger just by a constant factor from the uniform process,
we get with very high probability that $\w'_i > \w_i$, for every $i$,
where $\w_i$ is the number of samples taken from $D_i$.
This implies that if we prove a lower bound for algorithms that receive
samples generated by the Poissonized process, then we obtain a related lower
bound for algorithms that work in the uniform sampling model.

As opposed to the process that takes a fixed number of samples according
to the uniform sampling model, the benefit of the Poissonized process is
that the $\alpha_{i,j}$'s determined by this process are independent.
Therefore, the type of sample histogram that element $j$ has is completely
independent of the types of sample histograms the other elements have.
We get that the fingerprint distribution is a generalized multinomial distribution,
which %
has been studied by Roos~\cite{roos}
(the connection is due to Valiant~\cite{val}).




%
%
%
%
%
%
%
%
%
%
%
%
%
%

\def\DetailsWish{
\begin{definition} \label{def:pois}
%
%
%
%
%
%
%
%
In the \emph{Poissonized} uniform sampling model with parameter $\w$
(which we'll refer to as the {\em $\w$-Poissonized} model),
given a list $\calD = (D_1,\ldots,D_m)$ of $m$ distributions,
%
a sample is generated as follows.
\BI
\item Draw $m$ independent samples $\w_1,\ldots, \w_m$ from $\poi(\w)$,
\item Return $\w_i$ independent samples distributed according to $D_i$ for each $i\in[m]$.
\EI
%
%
%
%
\end{definition}

%
%
%
%
%
%

\begin{lemma} \label{uniform.lem}
Assume %
there exists a tester $T$ in the uniform sampling model
for a property $\calP$ of lists of $m$ distributions,
that takes a sample of size $s = \w m$ where $\w \geq c$ for some
sufficiently large constant $c$, and works for every
$\eps \geq \eps_0$ where $\eps_0$ is a constant
(and whose success probability is at least $2/3$).
Then there exists a tester $T'$ for $\calP$ in the Poissonized
uniform sampling model with parameter $2\w$,
that works for every $\eps \geq \eps_0$ and whose success
probability is at least ${19}/{30}$.
\end{lemma}
\BPF
Roughly speaking, the tester $T'$ tries to simulate $T$ if it has
a sufficiently large sample, and otherwise it guesses the answer.
More precisely,
%
%
%
%
%
%
%
%
%
%
%
%
%
%
let $T'$ be a tester in the Poissonized uniform sampling model with parameter $2\w$.
The tester $T'$ receives a set of samples, $S'$, where each sample in the set is
distributed uniformly and independently over the $m$ distributions, and their
number, $\w'$, is distributed as $\w' \sim \poi(2\w m)$ (see~\cite[p. 216]{feller}).
By \expref{Lemma}{l:poiexp} we have that,
\BEQ
\Pr\left[\w' < \w m\right] \leq (3/4)^{\w m/2}\,. \label{kappa.eq}
\EEQ
%
%
If $\w' \geq \w m$  then $T'$ picks a random subset of samples, $S\subset S'$, of size $\w m$ and simulates $T$ on $S$.
%
Otherwise it ``accepts'' or ``rejects'' with equal probability. %
Since the samples in $S$ are distributed as $\w m$ samples drawn according to the uniform sampling model, if equation~(\ref{kappa.eq}) holds then $T'$ simulates a tester in the uniform sampling model successfully.

The probability
that $\w' \geq \w m$ is at least  $1 - (3/4)^{\w m/2}$, which is
greater than ${4}/{5}$ for $\w > c$ and a sufficiently large constant $c$.
%
%
%
%
%
%
%
%
%
Therefore, the success probability of $T'$ is at least
\[
\frac{4}{5}\cdot\frac{2}{3} + \frac{1}{5}\cdot\frac{1}{2} =
\frac{19}{30}\,,
\]
as desired.
\EPF

\medskip
Given \expref{Lemma}{uniform.lem} it suffices to consider samples that are
%
generated in the Poissonized uniform sampling model. The process for generating a
sample $\{\alpha_{1,j}, \ldots, \alpha_{m,j}\}_{j\in[n]}$ (recall that $\alpha_{i,j}$ is the number of times that
element $j$ was selected by distribution $D_i$) in the $\w$-Poissonized model
is equivalent
%
to the following process: For each $i \in [m]$
and $j \in [n]$, independently select $\alpha_{i,j}$ according
to $\poi(\w\cdot D_i(j))$ (see~\cite[p. 216]{feller}). Thus the probability
of getting a particular histogram $\vec{a} = (a_{1},\ldots,a_{m})$
for element $j$ is $p^{\calD,\w}(j;\vec{a})$ (as defined in
equation~(\ref{calDi-def.eq})). We can represent the event that
the histogram of element $j$ is $\vec{a}$ by
a Bernoulli random vector $\vec{b}_j$ that is indexed by all $\vec{v}\in \mathbb{N}^m$, is $1$
in the coordinate corresponding to $\vec{a}$, and is $0$ elsewhere.
Given this representation, the fingerprint of the sample corresponds
to $\sum_{j=1}^n \vec{b}_j$.

As defined above, the dimension of the $\vec{b}_j$'s is unbounded.
This is simply because
$\poi(\lambda)$ is unbounded for every $\lambda > 0$, and so every
possible histogram has non-zero probability.
However, we may consider
vectors of finite dimension (we show in the proof of \expref{Theorem}{wish.thm}
that it is sufficient) and map the histograms that we do not consider to
the vector $(0,\ldots, 0)$.
Roos's theorem, stated next, shows
that the distribution of the fingerprints can be approximated by a multivariate
Poisson distribution (the Poisson here is related to the fact that the fingerprints' distributions are generalized multinomial distributions and not related to the Poisson from the Poissonization process). For simplicity, the theorem is stated for
vectors $\vec{b}_j$ that are indexed directly, that is
$\vec{b}_j = (b_{j,1},\ldots,b_{j,h})$.
%
\begin{theorem}[\cite{roos}] \label{roos.thm}
Let $D^{S_n}$ be the distribution of the sum $S_n$ of $n$ independent
Bernoulli random vectors $\vec{b}_1, \ldots, \vec{b}_n$ in $\mathbb{R}^h$
where $\Pr\bigl[\vec{b}_j = \vec{e}_\ell\bigr] = p_{j,\ell}$ and
$\Pr\bigl[\vec{b}_j = (0,\ldots, 0)\bigr] = 1 - \sum_{\ell=1}^h p_{j,\ell}$
(here $\vec{e}_\ell$ satisfies $e_{j,\ell} = 1$ and $e_{j,\ell'}=0$ for
every $\ell' \neq \ell$).
%
Let us  %
define an $h$-dimensional vector
$\vec{\lambda} = (\lambda_1,\ldots,\lambda_h)$
%
%
by %
 $\lambda_\ell := \sum_{j=1}^n p_{j,\ell}$.
Then
\BEQ\label{DSn-poi.eq}
\left\|D^{S_n} - \poi(\vec{\lambda})\right\|_1
    \leq \frac{88}{5} \sum_{\ell=1}^h \frac{\sum_{j=1}^n p_{j,\ell}^2}
       {\sum_{j=1}^n p_{j,\ell}} \; .
\EEQ
\end{theorem}

We next show how to obtain a bound on sums of the form given in
equation~(\ref{DSn-poi.eq}) under appropriate conditions.
\begin{lemma}\label{eps.lem}
Given a list $\calD = (D_1,\ldots,D_m)$ of $m$ distributions over $[n]$
and a real number $0 < \delta \leq 1/2$ such that for all $i\in\left[m\right]$
and for all $j\in\left[n\right]$,  $D_i(j) \leq {\delta}/{(m\cdot \w)}$
for some integer $\w$, we have that
\BEQ\label{eps.eq}
\sum_{\vec{a}\in \mathbb{N}^m\setminus\vec{0}}
  \frac{\sum_{j=1}^n p^{\calD,\w}(j;\vec{a})^2}{\sum_{j=1}^n p^{\calD,\w}(j;\vec{a})} \leq 2\delta\,.
\EEQ
%
%
\end{lemma}
\BPF
%

\BEQN
\sum_{\vec{a}\in \mathbb{N}^m\setminus\vec{0}}
   \frac{\sum_{j=1}^n p^{\calD,\w}(j;\vec{a})^2}{\sum_{j=1}^n p^{\calD,\w}(j;\vec{a})}
&\leq& \sum_{\vec{a}\in \mathbb{N}^m\setminus\vec{0}} \max_j \left(p^{\calD, \w}(j;\vec{a})\right) \\
&=& \sum_{\vec{a}\in \mathbb{N}^m\setminus\vec{0}}
   \max_j \left(\prod_{i=1}^m \poi(a_i;\w\cdot D_i(j))\right) \\ \label{eq:poi1}
&\leq& \sum_{\vec{a}\in \mathbb{N}^m\setminus\vec{0}} \left(\frac{\delta}{m}\right)^{a_1+\ldots +a_m} \\
&\leq& \sum_{a=1}^{\infty} m^a \left(\frac{\delta}{m}\right)^a \\
&\leq& 2\delta\;,\label{eq:poi2}
\EEQN
where the inequality in equation~(\ref{eq:poi2}) holds for $\delta\leq 1/2$ and the inequality in equation~(\ref{eq:poi1}) follows from:













%
\begin{eqnarray}
\poi(a;\w\cdot D_i(j)) &=& \frac{\eee^{-\w\cdot D_i(j)}(\w\cdot D_i(j))^a}{a!} \\
&\leq& (\w\cdot D_i(j))^a \\
&\leq& \left(\frac{\delta}{m}\right)^a \;,
\end{eqnarray}
and the proof is completed.
\EPF

\begin{proof}[Proof of \expref{Theorem}{wish.thm}]
%
%
Let $D^{\calD^-}$ and $D^{\calD^+}$ denote the distribution of the fingerprints when taking samples in the $\w$-Poissonized model (see \expref{Definition}{def:pois}) from $\calD^-$ and $\calD^+$, respectively.
We first bound the $\ell_1$-distance between $D^{\calD^-}$ and $D^{\calD^+}$ and then prove the desired result on the uniform sampling model by applying \expref{Lemma}{uniform.lem}.
Let $A\subset \mathbb{N}^m$ be some finite subset of histograms such that for both $\calD^-$ and $\calD^+$, the probability (in the $\w$-Poissonized model) that at least one of the elements has a sample histogram $\vec{a}\notin A$ is smaller than ${1}/{120}$ (it is easy to verify, for example by Markov's inequality, that such a subset exists).
Denote by $\mathcal{F}$ the subset of fingerprints describing outcomes where for each element in $\domain$,  the sample histogram is in $A$.
Hence,
\[
\sum_{\bar{\mathcal{F}} \owns \vec{f}}\left|
  D^{\calD^-}\left(\vec{f}\right) - D^{\calD^+}\left(\vec{f}\right)
\right| \leq \frac{1}{60}\,.
\]
We next turn to bound 
\[
\sum_{\mathcal{F} \owns \vec{f}}\left| D^{\calD^-}\left(\vec{f}\right)
  - D^{\calD^+}\left(\vec{f}\right) \right|\,.
\]

By the first premise of the theorem, $D^+_i(j),D^+_i(j) \leq  {\delta}/{(\w m)-}$
for every $i\in[m]$ and $j \in [n]$.
By \expref{Lemma}{eps.lem} this implies that
equation~(\ref{eps.eq}) holds both for $\calD = \calD^+$ and for
$\calD = \calD^-$.
Combining this with \expref{Theorem}{roos.thm}
we get that the $\ell_1$-distance between the
fingerprint distribution over $\mathcal{F}$ when the sample
is generated according to $\calD^+$ (in the $\w$-Poissonized model) and the
distribution $\poi\bigl(\vec{\lambda}^{\calD^+,\w}\bigr)$
is at most
\[
\frac{88}{5}\cdot 2\delta = \frac{176}{5}\delta\,,
\]
and an analogous statement holds for $\calD^-$.
By applying the premise in equation~(\ref{poi-poi.eq})
(concerning the $\ell_1$ distance between $\poi\bigl(\vec{\lambda}^{\calD^+,\w}\bigr)$
and $\poi\bigl(\vec{\lambda}^{\calD^-,\w}\bigr)$)
and the triangle inequality, we get that
%
%
%
%
%
%
%
%
\[
\sum_{\bar{\mathcal{F}}\cup \mathcal{F} \owns \vec{f}}\left|
  D^{\calD^-}\left(\vec{f}\right) - D^{\calD^+}\left(\vec{f}\right)
\right| \leq 2 \cdot \frac{176}{5}\delta + \frac{31}{60} -
\frac{352\delta}{5} + \frac{1}{60}= \frac{16}{30}\,,
\]
which implies that the statistical distance is smaller than ${8}/{30}$.
Thus a tester that accepts $\calD^+$ and rejects $\calD^-$ with
probability greater than
\[
\frac{1}{2} + \frac{1}{2} \cdot \frac{8}{30}= \frac{19}{30}\,,
\]
viewing only the above-mentioned fingerprints, does not exist.
By \expref{Lemma}{uniform.lem} the theorem follows.
\end{proof}
%
%
%
}
%

%

\DetailsWish

%


\subsection{Proof of \expref{Theorem}{lb1.thm}}\label{id-lb-proof.subsec}
%
In this subsection we show how to
apply \expref{Theorem}{wish.thm} to two lists of distributions,
$\calD^+$ and $\calD^-$,
which we will define shortly,
where $\calD^+ \in \calPid = \calPidmn$  while
$\calD^-$ is $(1/20)$-far from $\calPid$.
%
Recall that by the premise of \expref{Theorem}{lb1.thm},
%
$n \geq cm\log m$ for some sufficiently large constant $c > 1$.
In the proof it will be convenient to assume that $m$ is even and
that $n$ (which corresponds in \expref{Lemma}{NOT.lem} to $2t$) is divisible by 4. It is not hard to verify that it is
possible to reduce the general case to this case.
%
%
%
In order to define $\calD^-$, we shall need the next lemma.
%
%
\begin{lemma} \label{NOT.lem}
For every two even integers $m$ and $t$,
%
there exists a $0/1$-valued matrix $M$ with $m$ rows and $t$ columns for which the
following holds:
\BE
\item In each row and each column of $M$, exactly half of the elements are $1$ and
the other half are $0$.
\item\sloppy For every integer $2 \leq x < m/2$,
 and for every subset $S \subseteq [m]$ of size $x$,
%
the number of columns $j$ such that $M[i,j] =1$ for every $i \in S$
is at least
\[
t\cdot\left(\frac{1}{2^x}\left(1- \frac{2x^2}{m}\right) -
  \sqrt{\frac{2x\ln m}{t}}\right)\,,
\]
and at most
%
\[
t\cdot\left(\frac{1}{2^x} + \sqrt{\frac{2x\ln m}{t}}\right)\,.
\]
\EE
\end{lemma}
























%
\def\NOTLemPrf{



Consider

%
selecting a matrix $M$ randomly as follows.
Denote the first $t/2$ columns of $M$ by $F$. For each column in $F$, pick,
independently from the other $t/2-1$ columns in $F$, a random half of its elements to be $1$,
and the other half of the elements to be $0$.
%
Columns $t/2+1,\ldots,t$ are the negations of columns $1,\ldots,t/2$,
respectively. Thus, in each row and each column of $M$, exactly half
of the elements are $1$ and the other half are $0$.

Consider a fixed choice of $x$.
For each column $j$ between $1$ and $t$, each subset of rows $S \subseteq [m]$ of size $x$,
%
and $b \in \bitset$,
define the indicator random variable $I_{S,j,b}$ to be $1$ if and only if
%
%
$M[i,j] = b$ for every $i \in S$.
Hence,
%

\begin{eqnarray}
\Pr[I_{S,j,b}=1]
= \frac{1}{2} \cdot \left(\frac{1}{2} - \frac{1}{m}\right)\cdot \ldots
     \cdot \left(\frac{1}{2} - \frac{x-1}{m}\right)\,.
\end{eqnarray}







%
Clearly, $\Pr[I_{S,j,b}=1] < {1}/{2^{x}}$. On the other hand,
\begin{eqnarray}
\Pr[I_{S,j,b}=1]
&\geq& \left(\frac{1}{2} - \frac{x}{m}\right)^x \\
&=& \frac{1}{2^x}\left(1 - \frac{2x}{m}\right)^x \\
&\geq& \frac{1}{2^x}\left(1 - \frac{2x^2}{m}\right) \;,
\end{eqnarray}
where the last inequality is due to Bernoulli's inequality which states that $(1+x)^n > 1 + nx$, for every real, nonzero number $x>-1$ and an integer $n>1$ (\cite{bern}).
%

Let $E_{S,b}$ denote the expected value of $\sum_{j=1}^{t/2} I_{S,j,b}$.
{}From the fact that columns $t/2+1,\ldots,t$ are the negations of columns $1,\ldots,t/2$ it follows that
\[
\sum_{j=t/2+1}^{t} I_{S,j,1} = \sum_{j=1}^{t/2} I_{S,j,0}\,.
\]
Therefore, the expected number of columns %
$1\leq j\leq t$ such that $M[i,j] = 1$ for every $i \in S$ is
%
simply $E_{S,1} + E_{S,0}$ (that is, at most $t\cdot {1}/{2^{x}}$
and at least $(t\cdot {1}/{2^x})\left(1 - {2x^2}/{m}\right)$).
By the additive Chernoff bound,
%
%
%
%
%
%
%
%
%
%
%

\BEQN
\Pr\Bigg[\Big|\sum_{j=1}^{t/2} I_{S,j,b} - E_{S,b}\Big|
             > \sqrt{\frac{t x\ln m}{2}}\Bigg]
&<& 2\exp(-2(t/2) (2x\ln m)/t)\\
&=& 2m^{-2x}\,.
\EEQN








%
Thus, by taking a union bound (over $b \in \bitset$),
%
%
%






\BEQN
\Pr\Bigg[\Big|\sum_{j=1}^{t} I_{S,j,1} - (E_{S,1}+E_{S,0})\Big| >  \sqrt{2t x\ln m}\Bigg]
\;<\; 4m^{-2x}\,.
\EEQN

By taking a union bound over all subsets $S$
we get that $M$ has the desired properties with probability
greater than $0$.
} %


 \BPF
 \NOTLemPrf
 \EPF


\medskip
We first define $\calD^+$, in which all distributions are identical. Specifically,
for each $i \in [m]$:
\begin{equation}
D_i^+(j) \;\eqdef\;
\begin{cases}
 \frac{1}{n^{2/3}m^{1/3}} & \text{if } 1 \leq j \leq \frac{n^{2/3}m^{1/3}}{2}\,,\\[1ex]
 \frac{1}{n} & \text{if } \frac{n}{2} < j \leq n\,,\\[1ex]  0 & \text{otherwise.}
\end{cases}
\label{Dpl-def.eq}
\end{equation}
We now turn to defining $\calD^-$.
Let $M$ be a matrix as in \expref{Lemma}{NOT.lem} for $t=n/2$.
For every $i \in [m]$:
%
\BEQ
D_i^-(j) \;\eqdef\;
\begin{cases}  \frac{1}{n^{2/3}m^{1/3}} & \text{if } 1 \leq j \leq
        \frac{n^{2/3}m^{1/3}}{2}\,,\\[1ex]
       \frac{2}{n} & \text{if } \frac{n}{2} < j \leq n \text{ and }M[i,j-n/2]=1\,,\\[1ex]  0 & \text{otherwise.}\end{cases}
\label{Dmi-def.eq}
\EEQ
%
%
For both $\calD^+$ and $\calD^-$, we refer to the elements
$1 \leq j \leq {n^{2/3}m^{1/3}}/{2}$ as the \emph{heavy} elements,
and to the elements ${n}/{2} \leq j \leq n$, as the \emph{light} elements.
Observe that each heavy element has exactly the same probability weight,
${1}/{(n^{2/3}m^{1/3})}$, in all distributions $D_i^+$ and $D_i^-$.
On the other hand, for each light element $j$, while $D_i^+(j) = {1}/{n}$
(for every $i$), in $\calD^-$ we have that $D_i^+(j) = {2}/{n}$
for half of the distributions, the distributions selected by the $M$, and $D_i^+(j) = 0$ for half of the distributions, the distributions which are not selected by $M$. We later use the properties of $M$ to bound the $\ell_1$ distance between the fingerprints' distributions of $\calD^+$ and $\calD^-$.

%

\paragraph{A high-level discussion}
To gain some intuition before delving into the detailed proof,
consider first the special case that $m=2$
(which was studied by Valiant~\cite{val}, and indeed the construction is
the same as the one he analyzes (and was initially proposed
in~\cite{BFRSW})).
%
%
In this case
each heavy element has probability weight $\Theta(1/n^{2/3})$ and we would like to
establish a lower bound of $\Omega(n^{2/3})$ on the number of samples required
to distinguish between $\calD^+$ and $\calD^-$. That is, we would like to show that
the corresponding fingerprints' distributions when the sample is of size
$o(n^{2/3})$ are very similar.

The first main observation is that since the probability
weight of light elements is $\Theta(1/n)$ in both $\calD^+$ and $\calD^-$,
the probability that a light element will appear more than twice  in a sample of
size $o(n^{2/3})$ is very small. That is (using the fingerprints of histograms
notation we introduced previously),
 for each $\vec{a} = (a_1,a_2)$ such that $a_1+a_2 > 2$, the sample
will not include (with high probability) any light element
$j$ such that $\alpha_{1,j} = a_1$ and $\alpha_{2,j} = a_2$
(for both $\calD^+$ and $\calD^-$).
Moreover, for every $x\in \{1,2\}$, the expected number of elements $j$ such that $(\alpha_{1,j},\alpha_{2,j}) = (x,0)$ is the same in $\calD^+$ and $\calD^-$, as well as the variance (from symmetry, the same applies to $(0,x)$).
Thus, most of the difference between the
fingerprints' distributions is due to the numbers of elements $j$
such that $(\alpha_{1,j},\alpha_{2,j}) = (1,1)$.
For this setting we do expect to see a non-negligble difference
for light elements between $\calD^+$ and $\calD^-$ (in particular, we cannot
get the $(1,1)$ histogram for light elements in $\calD^-$, as opposed to
$\calD^+$).
%

Here is where the heavy elements come into play. Recall that in both $\calD^+$ and $\calD^-$
the heavy elements have the same probability weight, so that the expected
number of heavy elements $i$ such that $(a_{1,j},a_{2,j}) = (1,1)$ is the same for $\calD^+$ and $\calD^-$. However,
intuitively, the variance of  these numbers for the heavy elements
``swamps'' the differences between the light elements so that it is not possible
to distinguish between  $\calD^+$ and $\calD^-$. The actual proof, which formalizes
(and quantifies)
this intuition, considers the difference between the values of the vectors $\vec{\lambda}^{\calD^+,k}$ and
$\vec{\lambda}^{\calD^-,k}$ (as defined in
equation~(\ref{vecl-calD-def.eq}))
%
in the coordinates corresponding to
$\vec{a}$ such that $a_1 + a_2 = 2$.
We can then apply Lemmas~\ref{disval.lem} and~\ref{dis.lem} to
obtain equation~(\ref{poi-poi.eq}) in \expref{Theorem}{wish.thm}.

Turning to $m> 2$, it is no longer true that in a sample of size $o(n^{2/3} m^{1/3})$
we will not get histogram vectors $\vec{a}$ such that $\sum_{i=1}^m a_i > 2$
for light elements. Thus we have to deal with many more vectors $\vec{a}$
(of dimension $m$)
and to bound the total contribution of all of them to the difference between
fingerprints of $\calD^+$ and of $\calD^-$.
To this end we partition the set of all possible histograms'
vectors into several subsets
according to their Hamming weight $\sum_{i=1}^m a_i$ and depending
on whether all $a_i's$ are in $\bitset$,
%
or there exists a least one $a_i$ such that $a_i \geq 2$.
In particular, to deal with the former (whose number, for each choice of
Hamming weight $x$ is relatively large, \ie, roughly $m^x$), we use
the properties of the matrix $M$ based on which $\calD^-$ is defined.
We note that from the analysis we see that, similarly to when $m=2$, we need the variance of the heavy elements to play a role just for the cases where $\sum_{i=1}^m a_i = 2$ while in the other cases the total contribution of the light elements is rather small.

\noindent
In the remainder of this section we provide the details of the analysis.

\medskip

%

Before establishing that indeed
$\calD^-$ is $\Omega(1)$-far from $\calPid$,
we introduce some more notation (which will be used throughout
the remainder of the proof of \expref{Theorem}{lb1.thm}).




%
Let $S_x$ be the set of vectors of dimension $m$ that contain exactly $x$ coordinates that are $1$,
and all the rest are $0$ (which corresponds to an element that was sampled once or $0$ times by each distribution). Let $A_x$ be the set of vector of dimension $m$ that their coordinates sum up to $x$ but must contain at least one
coordinate that is $2$ (which corresponds to an element that was samples at least twice by at least one distribution).
More formally, for any integer $x$, we define the following two subsets
of $\mathbb{N}^m$:
\BEQ
S_x \eqdef \left\{\vec{a}\in \mathbb{N}^m\; :\;
\begin{array}{rl} \sum_{i=1}^m a_i = x \text{ and } \\ \forall i\in[m], a_i < 2 \end{array}
\right\} \;,
\label{Sx-def.eq}
\EEQ
and
\BEQ
A_x \eqdef \left\{\vec{a}\in \mathbb{N}^m \;:\;
  \begin{array}{rl} \sum_{i=1}^m a_i = x \text{ and } \\ \exists i\in[m], a_i\geq 2  \end{array} \right\}\,.
\label{Ax-def.eq}
\EEQ
For $\vec{a}\in \mathbb{N}^m$,
let ${\rm sup}(\vec{a}) \eqdef \{i:a_i \neq 0\}$ denote the \emph{support}
of $\vec{a}$, and let
%
%
%
%
%

\BEQ
I_M(\vec{a}) \eqdef \left\{j\,:\; D^-_i(j) = \frac{2}{n}
   \;\;\;\forall i \in {\rm sup}(\vec{a})\right\}\,.
\label{I-veca-def.eq}
\EEQ




Note that in terms of the matrix $M$ (based on which $\calD^-$ is defined),
$I_M(\vec{a})$ consists of the columns in $M$ whose restriction to the support
of $\vec{a}$ contains only $1$'s. In terms of the $\calD^-$, it corresponds to the set of light elements that might have a sample histogram of $\vec{a}$ (when sampling according to $\calD^-$).

%







%

%

\begin{lemma} \label{far.lem}
For every $m > 5$ and for
%
$n\geq c\ln m$ for some sufficiently large $c$,
we have that 
\[
\sum_{i=1}^m \|D^-_i - D^*\|_1 > \frac{m}{20}
\]
for every
distribution $D^*$ over $[n]$.
That is, the list $\calD^-$ is $(1/20)$-far from $\calPid$.
\end{lemma}

%

%
\def\FarLemPrf{
Consider any $\vec{a}\in S_2$. By \expref{Lemma}{NOT.lem}, setting $t=n/2$,
the size of $I_M(\vec{a})$, \ie, the number of light elements $\ell$
such that $D^-_i[\ell] = {2}/{n}$ for every $i\in {\rm sup}(\vec{a})$, is at most
%
\[
\frac{n}{2}\left(\frac{1}{4} + \sqrt{\frac{8 \ln m}{n}}\right)\,.
\]
The same upper bound
holds for the number of light elements $\ell$
such that $D^-_{i}[\ell] = 0$ for every $i\in {\rm sup}(\vec{a})$.
This implies that for every $i\neq i'$ in $[m]$,
for at least
%
\[
\frac{n}{2} - n\left(\frac{1}{4} + \sqrt{\frac{8 \ln m}{n}}\right)
\]
of the light elements, $\ell$,
we have that $D^-_i[\ell] = {2}/{n}$ while $D^-_{i'}[\ell] = 0$,
or that $D^-_{i'}[\ell] = {2}/{n}$ while $D^-_{i}[\ell] = 0$.
Therefore,
%
%
%
\[
\|D^-_i - D^-_{i'}\|_1 \geq \frac{1}{2} -2\sqrt{\frac{8\ln m}{n}}\,,
\]
 which for $n \geq c\ln m$ and a sufficiently large constant $c$,
is at least ${1}/{8}$. Thus, by the triangle inequality we have that
for every $D^*$,
\[
\sum_{i=1}^m \|D^-_i - D^*\|_1 \geq \left\lfloor{\frac{m}{2}}\right\rfloor\cdot
\frac{1}{8}\,,
\] which greater than $m/20$ for $m>5$.
}
%

\BPF
\FarLemPrf
\EPF
\medskip

%

In what follows we work towards establishing that
 equation~(\ref{poi-poi.eq}) in \expref{Theorem}{wish.thm} holds for
$\calD^+$ and $\calD^-$.
Set $\w = \delta \cdot {n^{2/3}}/{m^{2/3}}$, where $\delta$ is a constant to be determined later.
We shall use the shorthand
$\vec{\lambda}^+$ for $\vec{\lambda}^{\calD^+,\w}$, and
$\vec{\lambda}^-$ for $\vec{\lambda}^{\calD^-,\w}$
(recall that the notation $\vec{\lambda}^{\calD,\w}$
was introduced in equation~(\ref{vecl-calD-def.eq})).
By the definition of $\vec{\lambda}^+$, for each
$\vec{a} \in \mathbb{N}^m$,

\BEQN
\vec{\lambda}^+(\vec{a})
&=&
   \sum_{j=1}^n\prod_{i=1}^m \frac{(\w\cdot D^+_i(j))^{a_i}}{\eee^{\w\cdot D^+_i(j)} \cdot a_i!}\\
&=& \sum_{j=1}^{n^{2/3}m^{1/3}/2}\prod_{i=1}^m  \frac{(\delta/m)^{a_i}}{\eee^{\delta/m}\cdot a_i!}
\;+\; \sum_{j=n/2+1}^n \prod_{i=1}^m \frac{(\delta/(n^{1/3}m^{2/3}))^{a_i}}{\eee^{\delta/(n^{1/3} m^{2/3})}\cdot a_i!} \\
&=&\frac{n^{2/3} m^{1/3}}{2\eee^{\delta}}\prod_{i=1}^m \frac{(\delta/m)^{a_i}}{a_i!}
\label{vec-lambda-pl.eq}
\;+\; \frac{n}{2\eee^{\delta(m/n)^{1/3}}} \prod_{i=1}^m \frac{(\delta/(n^{1/3}m^{2/3}))^{a_i}}{a_i!} \,.
\EEQN














%
By the construction of $M$, for every light $j$,
%
$\sum_{i=1}^m D^-_i(j) = \frac{2}{n}\cdot\frac{m}{2} = \frac{m}{n}$.
%
Therefore,
%
%
%
%
%
%
%
%

\BEQN
\vec{\lambda}^-(\vec{a})
&=& \frac{n^{2/3} m^{1/3}}{2\eee^{\delta}}
  \prod_{i=1}^m \frac{(\delta/m)^{a_i}}{a_i!}
\label{vec-lambda-mi.eq}
\;+\; \frac{1}{\eee^{\delta(m/n)^{1/3}}} \sum_{j\in I_M(\vec{a})} \prod_{i=1}^m
    \frac{(2\delta/(n^{1/3}m^{2/3}))^{a_i}}{a_i!}   \,.
\EEQN












%
Hence, $\vec{\lambda}^+(\vec{a})$ and $\vec{\lambda}^-(\vec{a})$ differ only
on the term which corresponds to the contribution of the light elements.

Equations~(\ref{vec-lambda-pl.eq}) and~(\ref{vec-lambda-mi.eq}) demonstrate why we choose $M$
with the specific properties defined in \expref{Lemma}{NOT.lem}.




First of all, in order for every $D^-_i$ to be a probability distribution, we want
each row of $M$ to sum up to exactly $n/4$.
%
%
We also want each column of $M$ to sum up to exactly $m/2$,
in order to get $\prod_{i=1}^m \eee^{-\w\cdot D^+_i(j)} = \prod_{i=1}^m \eee^{-\w\cdot D^-_i(j)}$.
%
Finally,
we would have liked $\left|I_M(\vec{a})\right| \cdot \prod_{i=1}^m 2^{a_i}$ to
equal $n/2$ for every $\vec{a}$. This would imply that  $\vec{\lambda}^+(\vec{a})$
and $\vec{\lambda}^-(\vec{a})$  are equal.
As we show below, this is in fact true for every $\vec{a} \in S_1$.
For vectors $\vec{a} \in S_x$ where $x > 1$, the second
condition in \expref{Lemma}{NOT.lem} ensures that
$\left|I_M(\vec{a})\right|$ is sufficiently close to  $({n}/{2})\cdot ({1}/{2^x})$.
%
%
%
%
 This property of $M$ is not necessary in order to bound the contribution
 of the vectors in $A_x$. The bound that we give for those vectors is less tight,
but since there are fewer such vectors, it suffices.

\smallskip
We start by considering the contribution to equation~(\ref{poi-poi.eq})  of
histogram vectors $\vec{a} \in S_1$
(\ie, vectors of the form $(0,\ldots,0,1,0,\ldots,0)$) which correspond to
the number of elements that are sampled only by one distribution, once.
%
%
%

We prove that in the Poissonized uniform sampling model,



for every $\vec{a} \in S_1$ the number of elements with such sample
histogram is distributed exactly the same in $\calD^+$ and $\calD^-$.

\begin{lemma} \label{s1a2.lem}
\BEQN
\sum_{\vec{a} \in S_1}
  \left\|\poi(\vec{\lambda}^+(\vec{a}) - \poi(\vec{\lambda}^-(\vec{a})) \right\|_1 = 0\,.
\EEQN
\end{lemma}
\BPF
For every $\vec{a}  \in S_1$, the size of
$I_M(\vec{a})$ is ${n}/{4}$, thus,
%

\BEQN
\sum_{j\in I_M(\vec{a})}
   \prod_{i=1}^m \frac{(2\delta/(n^{1/3}m^{2/3}))^{a_i}}{a_i!}
  = \frac{n}{2} \prod_{i=1}^m \frac{(\delta/(n^{1/3}m^{2/3}))^{a_i}}{a_i!} \; .
\EEQN







%
%

By equations~(\ref{vec-lambda-pl.eq}) and~(\ref{vec-lambda-mi.eq}),
it follows that
$\left|\vec{\lambda}^+(\vec{a}) - \vec{\lambda}^-(\vec{a})\right| = 0$
for every $\vec{a}  \in S_1$.






%
The lemma follows by  applying equation~(\ref{poi-stat.eq}).
%
\EPF

\medskip\noindent
We now turn to bounding the contribution to equation~(\ref{poi-poi.eq})  of
histogram vectors $\vec{a} \in A_2$
(\ie, vectors of the form $(0,\ldots,0,2,0,\ldots,0)$ which correspond to
the number of elements that are sampled only by one distribution, twice.
\begin{lemma} \label{a2.lem}
\BEQN
  \left\|\poi(\vec{\lambda}^+(A_2)) - \poi(\vec{\lambda}^-(A_2)) \right\|_1 \leq 3\delta\,.
\EEQN
\end{lemma}
\BPF
For every $\vec{a}  \in A_2$, the size of
$I_M(\vec{a})$ is ${n}/{4}$, thus,
%

\BEQN
\sum_{j\in I_M(\vec{a})}
   \prod_{i=1}^m \frac{(2\delta/(n^{1/3}m^{2/3}))^{a_i}}{a_i!}
  = n \prod_{i=1}^m \frac{(\delta/(n^{1/3}m^{2/3}))^{a_i}}{a_i!} \; .\label{a2.eq1}
\EEQN
By equations~(\ref{vec-lambda-pl.eq}),~(\ref{vec-lambda-mi.eq}) and~(\ref{a2.eq1})
it follows that









%
%

\BEQN
\vec{\lambda}^-(\vec{a}) - \vec{\lambda}^+(\vec{a})
= \frac{n}{2\eee^{\delta(m/n)^{1/3}}} \prod_{i=1}^m
          \frac{(\delta/(n^{1/3}m^{2/3}))^{a_i}}{a_i!}
= \frac{n^{1/3}\delta^2}{4\eee^{\delta(m/n)^{1/3}}m^{4/3}}
  \;, \label{a2.eq2}
\EEQN









%
and that
\BEQN
\vec{\lambda}^-(\vec{a}) \geq \frac{n^{2/3} m^{1/3}}{2\eee^{\delta}}
  \prod_{i=1}^m \frac{(\delta/m)^{a_i}}{a_i!} =  \frac{n^{2/3}\delta^2}{4\eee^{\delta}m^{5/3}} \,.\label{a2.eq3}
\EEQN
%

By equations~(\ref{a2.eq2}) and~(\ref{a2.eq3}) we have that



%
\BEQN
\frac{\left(\vec{\lambda}^-(\vec{a}) - \vec{\lambda}^+(\vec{a})\right)^2}{\vec{\lambda}^-(\vec{a})} \leq
\frac{\eee^{\delta-2\delta(m/n)^{1/3}}\delta^2}{4m} 
\leq \frac{\delta^2}{m} \,. \label{a2.eq4}
\EEQN
By equation~(\ref{a2.eq4}) and the fact that $|A_2| = m$ we get
\BEQ
\sum_{\vec{a}\in A_2} \frac{\left(\vec{\lambda}^-(\vec{a}) - \vec{\lambda}^+(\vec{a})\right)^2}{\vec{\lambda}^-(\vec{a})} \leq
m \cdot \frac{\delta^2}{m} = \delta^2
\EEQ
The lemma follows by  applying \expref{Lemma}{dis.lem}.
%
\EPF

\smallskip
Recall that for
%
%
%
%
%
%
%
a subset $I$ of $\mathbb{N}^m$, $\poi(\vec{\lambda}(I))$ denotes
the multivariate Poisson distributions restricted to the coordinates of $\vec{\lambda}$
that are indexed by the vectors in $I$.
We separately deal with $S_x$ where $2 \leq x <m/2$, and $x \geq m/2$, where
our main efforts are with respect to the former, as the latter correspond
to very low probability events.

\begin{lemma} \label{Sx.lem}
For $m \geq 16$, $n \geq cm\ln m$ (where $c$ is a sufficiently large
constant) and for $\delta \leq 1/16$
\BEQ
\left\|\poi\left(\vec{\lambda}^+\left(\bigcup_{x=2}^{m/2-1}S_x\right)\right) -
    \poi\left(\vec{\lambda}^-\left(\bigcup_{x=2}^{m/2-1} S_x\right)\right) \right\|_1 \leq 32\delta\,.
\EEQ
\end{lemma}

\BPF
Let $\vec{a}$ be a vector in $S_x$ then by the definition of $S_x$, every coordinate of $\vec{a}$ is $0$ or $1$.

Therefore we make the following simplification of equation~(\ref{vec-lambda-pl.eq}):



For each $\vec{a}\in \bigcup_{x = 2}^{m/2-1}S_x$,

\BEQN
\vec{\lambda}^+(\vec{a}) =
  \frac{n^{2/3}m^{1/3}}{2\eee^{\delta}}\cdot\left(\frac{\delta}{m}\right)^x
     \;+\; \frac{n}{2\eee^{\delta(m/n)^{1/3}}}\cdot\left(\frac{\delta}{n^{1/3}m^{2/3}}\right)^x\,.
\EEQN







By \expref{Lemma}{NOT.lem},
for every $\vec{a}\in \bigcup_{x=2}^{m/2-1}S_x$ the size of
$I_M(\vec{a})$ is at most
\[
\frac{n}{2}\cdot\left(\frac{1}{2^x} + \sqrt{\frac{4x\ln m}{n}}\right)
\]
%
%
and at least
\[
\frac{n}{2}\cdot \left(\frac{1}{2^x}- \frac{2x^2}{2^xm}
       - \sqrt{\frac{4x\ln m}{n}}\right)\,.
\]
By equation~(\ref{vec-lambda-mi.eq}) this implies that
%

\BEQN
\vec{\lambda}^-(\vec{a})
 = \frac{n^{2/3}m^{1/3}}{2 \eee^{\delta}}
     \cdot\left(\frac{\delta}{m}\right)^x
    \;+\; \frac{n}{2\eee^{\delta(m/n)^{1/3}}}\cdot\left(\frac{1}{2^x}
          + \eta \right)\left(\frac{2\delta}{n^{1/3}m^{2/3}}\right)^x\;,
\EEQN









%
where
\[
- \left(\frac{2x^2}{2^xm} + \sqrt{\frac{4x\ln m}{n}}\right) \leq  \eta
\leq \sqrt{\frac{4x\ln m}{n}}
\]
 and thus 
\[
|\eta| \leq \sqrt{\frac{x}{m}} \cdot \left(\frac{2x^2}{2^x\sqrt{m}} +
  \sqrt{\frac{4m\ln m}{n}}\right)\,.
\]
By the facts that
$n \geq cm\ln m$ for some sufficiently large constant $c$, and that
\[
\frac{2x^2}{2^x\sqrt{m}}\leq \frac{1}{2}
\]
for every $2 \leq x < m/2$ and $m\geq 16$, we obtain that $|\eta| \leq \sqrt{\frac{x}{m}}$.
So we have that
%

\begin{eqnarray}
(\vec{\lambda}^+(\vec{a}) - \vec{\lambda}^-(\vec{a}))^2
\leq
 \left(\frac{n}{2\eee^{\delta(m/n)^{1/3}}}
    \cdot\left(\frac{2\delta}{n^{1/3}m^{2/3}}\right)^x\cdot \sqrt{\frac{x}{m}}\right)^2 
\leq
    \frac{n^2}{4}\cdot \left(\frac{4\delta^2}{n^{2/3}m^{4/3}}\right)^x \cdot \frac{x}{m}\,.
\end{eqnarray}










%
and that
\BEQN
\vec{\lambda}^-(\vec{a}) \geq \frac{n^{2/3}m^{1/3}}{2\eee^{\delta}}\cdot\left(\frac{\delta}{m}\right)^x\,.
\EEQN
Then we get, for $\delta \leq 1/2$, that
%

\BEQN
\frac{\left(\vec{\lambda}^+(\vec{a}) - \vec{\lambda}^-(\vec{a})\right)^2}
       {\vec{\lambda}^-(\vec{a})}
&\leq&
  \frac{\eee^{\delta}n^{4/3}}{2m^{1/3}}\cdot
    \left(\frac{4\delta}{n^{2/3}m^{1/3}}\right)^x\cdot
       \frac{x}{m} \\
&\leq& \frac{n^{4/3}}{m^{1/3}}\cdot\left(\frac{4\delta}{n^{2/3}m^{1/3}}\right)^x
     \cdot\frac{x}{m} \\
&\leq& \frac{n^{4/3}}{m^{4/3}}\cdot\left(\frac{4x^{1/x}\delta}{n^{2/3}m^{1/3}}\right)^x\\
&\leq& \frac{n^{4/3}}{m^{4/3}}\cdot\left(\frac{8\delta}{n^{2/3}m^{1/3}}\right)^x\,.
\EEQN













%
Summing over all \[\vec{a}\in \bigcup_{x =2}^{m/2-1} S_x\] we get:
%

\BEQN
\sum_{\vec{a}\in \bigcup_{x=2}^{m/2-1} S_x}
   \frac{(\vec{\lambda}^-(\vec{a}) - \vec{\lambda}^+(\vec{a}))^2}{\vec{\lambda}^-(\vec{a})}
&\leq& \sum_{x=2}^\infty \frac{n^{4/3}}{m^{4/3}}\cdot\left(\frac{8\delta m^{2/3}}{n^{2/3}}\right)^x \\
&=& \sum_{x=0}^\infty 64\delta^2\cdot\left(\frac{8\delta m^{2/3}}{n^{2/3}}\right)^x
          \label{use-bounds2.eq}\\
&\leq& \frac{64\delta^2}{1-8\delta} \\ &\leq& 128\delta^2\,.\label{use-bounds3.eq}
\EEQN
where in equation~(\ref{use-bounds2.eq}) we used the fact that
$n > m$, and equation~(\ref{use-bounds3.eq}) holds for $\delta \leq 1/16$.











%
%
The lemma follows by applying \expref{Lemma}{dis.lem}.
\EPF

\begin{lemma} \label{Sxm2.lem}
For $n \geq m$, $m\geq 12$ and $\delta \leq 1/4$,
$\displaystyle
\sum_{x\geq m/2} \sum_{\vec{a} \in S_x}
  \left\|\poi(\vec{\lambda}^+(\vec{a}) )
     - \poi(\vec{\lambda}^-(\vec{a})) \right\|_1 \leq 32\delta^3\,.
$
\end{lemma}
\def\SxLemBigPrf{
We first observe that
%
$|S_x| \leq m^x/x$ for every $x\geq 6$.
To see why this is true, observe that $|S_x|$ equals the number
of possibilities of arranging $x$ balls in $m$ bins, \ie,

\begin{eqnarray}
|S_x| = {m\choose x}
\leq \left(\frac{\eee m}{x}\right)^x \leq \frac{m^x}{x} \label{eq:sxm2:1} \,.
\end{eqnarray}







where we have used the premise that $m\geq 12$ and thus $x\geq 6$.
%

By equations~(\ref{vec-lambda-pl.eq}) and~(\ref{vec-lambda-mi.eq})



%
(and the fact that $|x-y| \leq \max\{x,y\}$ for every positive real numbers $x$,$y$),
%

\BEQN
\sum_{x\geq m/2} \sum_{\vec{a} \in S_x}
%
%
    \left|\vec{\lambda}^+(\vec{a}) -
        \vec{\lambda}^-(\vec{a}) \right|
&\leq& \sum_{x\geq m/2} \sum_{\vec{a} \in S_x}
  \frac{n}{2} \prod_{i=1}^m \left(\frac{2\delta}{n^{1/3}m^{2/3}}\right)^{a_i} \\
&=& \sum_{x\geq m/2} \sum_{\vec{a} \in S_x}
     \frac{n}{2}\left(\frac{2\delta}{n^{1/3}m^{2/3}}\right)^{\sum_{i=1}^{m} a_i}\\
&\leq& \sum_{x = m/2}^\infty \frac{m^x}{x}\cdot \frac{n}{2}\left(\frac{2\delta}{n^{1/3}m^{2/3}}\right)^x\\
&\leq& \sum_{x = m/2}^\infty \frac{2m^x}{m}\cdot \frac{n}{2}\left(\frac{2\delta}{n^{1/3}m^{2/3}}\right)^x\\
&=& \frac{n}{m} \sum_{x = m/2}^\infty \left(\frac{2\delta m^{1/3}}{n^{1/3}}\right)^x\\
&=& 8\delta^3 \sum_{x = m/2-3}^\infty  \left(\frac{2\delta m^{1/3}}{n^{1/3}}\right)^x \\
&\leq& \frac{8\delta^3}{1-2\delta} \label{eq:sxm2:2}\\
&\leq& 16\delta^3 \label{eq:sxm2:3}
\EEQN
where in equation~(\ref{eq:sxm2:2}) we used the fact that $n\geq m$
 and equation~(\ref{eq:sxm2:3}) holds for $\delta \leq 1/4$.
%





















%
%
The lemma follows by applying equation~(\ref{poi-stat.eq}).
%
}


\BPF
\SxLemBigPrf
\EPF


\smallskip\noindent

We finally turn to the contribution of $\vec{a} \in A_x$ such that $x \geq 3$.




%
%
%
%
%
%
%
%
%
\begin{lemma} \label{Ax.lem}
For $n \geq m$ and  $\delta \leq 1/4$,
\BEQ
\sum_{x\geq 3} \sum_{\vec{a} \in A_x}
  \left\|\poi(\vec{\lambda}^+(\vec{a}) )
     - \poi(\vec{\lambda}^-(\vec{a})) \right\|_1 \leq 16\delta^3\,.
\EEQ
\end{lemma}
%
\def\AxLemPrf{
We first observe that  $|A_x| \leq m^{x-1}$ for every $x$.
To see why this is true, observe that $|A_x|$ equals the number
of possibilities of arranging $x-1$ balls, where one ball is a ``special''
(``double'') ball in $m$ bins.
%
%
%
%

By equations~(\ref{vec-lambda-pl.eq}) and~(\ref{vec-lambda-mi.eq})



%
%
(and the fact that $|x-y| \leq \max\{x,y\}$ for every positive real numbers $x$,$y$),
%

\BEQN
\sum_{x\geq 3} \sum_{\vec{a} \in A_x}
%
%
    \left|\vec{\lambda}^+(\vec{a}) -
        \vec{\lambda}^-(\vec{a}) \right|
   &\leq& \sum_{x\geq 3} \sum_{\vec{a} \in A_x}
  \frac{n}{2} \prod_{i=1}^m \left(\frac{2\delta}{n^{1/3}m^{2/3}}\right)^{a_i} \\
&=& \sum_{x\geq 3} \sum_{\vec{a} \in A_x}
     \frac{n}{2}\left(\frac{2\delta}{n^{1/3}m^{2/3}}\right)^{\sum_{i=1}^{m} a_i}\\
&\leq& \sum_{x = 3}^\infty m^{x-1}\cdot \frac{n}{2}\left(\frac{2\delta}{n^{1/3}m^{2/3}}\right)^x\\
&=& \frac{n}{2m} \sum_{x = 3}^\infty \left(\frac{2\delta m^{1/3}}{n^{1/3}}\right)^x\\
&=& 4\delta^3 \sum_{x = 0}^\infty  \left(\frac{2\delta m^{1/3}}{n^{1/3}}\right)^x \\
&\leq& \frac{4\delta^3}{1-2\delta} \label{use-bounds4.eq}\\
&\leq& 8\delta^3 \label{use-bounds5.eq}
\EEQN
where in equation~(\ref{use-bounds4.eq}) we used the fact that $n\geq m$
 and equation~(\ref{use-bounds5.eq}) holds for $\delta \leq 1/4$.




















%
The lemma follows by
applying equation~(\ref{poi-stat.eq}).
%
}
%

\BPF
\AxLemPrf
\EPF

%

\medskip\noindent
We are now ready to finalize the proof of \expref{Theorem}{lb1.thm}.

\begin{proof}[Proof of \expref{Theorem}{lb1.thm}]
Let $\calD^+$  and $\calD^-$
be as defined in equations~(\ref{Dpl-def.eq}) and~(\ref{Dmi-def.eq}), respectively,
and recall that $\w = \delta \cdot {n^{2/3}}/{m^{2/3}}$ (where $\delta$
will be set subsequently).
By the definition of the distributions in $\calD^+$ and $\calD^-$,
the probability weight assigned to each element is at most
${1}/{(n^{2/3}m^{1/3})} = {\delta}/{(\w\cdot m)}$, as required
by \expref{Theorem}{wish.thm}.
By \expref{Lemma}{far.lem}, $\calD^-$ is $(1/20)$-far from $\calPid$.
Therefore, it remains to establish
that equation~(\ref{poi-poi.eq}) holds for $\calD^+$ and $\calD^-$.
Consider the following partition of $\mathbb{N}^m$:
{\small
\BEQ
\left\{
\{\vec{a}\}_{\vec{a}\in S_1},
A_2,
 \bigcup_{x=2}^{m/2-1}S_x,
\{\vec{a}\}_{\vec{a}\in \bigcup_{x \geq m/2}S_x},
\{\vec{a}\}_{\vec{a}\in \bigcup_{x \geq 3}A_x}
\right\} \;,
\EEQ
}
where $\{\vec{a}\}_{\vec{a}\in T}$ denotes the list of all singletons of elements in $T$.
By \expref{Lemma}{disval.lem} it follows that
%

\BEQN
\left\|\poi(\vec{\lambda}^+) - \poi(\vec{\lambda}^-) \right\|_1
&\leq& \sum_{\vec{a}\in S_1}
  \left\|\poi(\vec{\lambda}^+(\vec{a})
     - \poi(\vec{\lambda}^-(\vec{a})) \right\|_1  \\
&&\;+\; \left\|\poi(\vec{\lambda}^+(A_2)
     - \poi(\vec{\lambda}^-(A_2)) \right\|_1  \\
&&\;+\; \Big\|\poi(\vec{\lambda}^+(\bigcup_{x=2}^{m/2-1}S_x)) -
     \poi(\vec{\lambda}^-(\bigcup_{x=2}^{m/2-1}S_x)) \Big\|_1 \\
     &&\;+\; \sum_{x\geq m/2} \sum_{\vec{a}\in S_x}
     \left\|\poi(\vec{\lambda}^+(\vec{a}) - \poi(\vec{\lambda}^-(\vec{a})) \right\|_1 \\
&&\;+\; \sum_{x\geq 3} \sum_{\vec{a}\in A_x}
     \left\|\poi(\vec{\lambda}^+(\vec{a}) - \poi(\vec{\lambda}^-(\vec{a})) \right\|_1\,.
\EEQN
















%
For $\delta < 1/16$ we get by
%
Lemmas~\ref{s1a2.lem}--\ref{Ax.lem} that
\BEQN
\left\|\poi(\vec{\lambda}^+) - \poi(\vec{\lambda}^-) \right\|_1 &\leq& 35\delta + 48\delta^3\;,
\EEQN
which is less than 
\[
\frac{31}{60} - \frac{352\delta}{5}
\]
for $\delta = 1/200$.
Hence, the theorem follows from \expref{Theorem}{wish.thm}.
\end{proof}

%
%
%
%
%
\providecommand{\bibhead}[1]{}
%
\expandafter\ifx\csname pdfbookmark\endcsname\relax%
  \providecommand{\tocrefpdfbookmark}{}
\else\providecommand{\tocrefpdfbookmark}{%
   \hypertarget{tocreferences}{}%
   \pdfbookmark[1]{References}{tocreferences}}%
\fi

\tocrefpdfbookmark
\begin{thebibliography}{10}

\bibitem{ACS10}\bibhead{ACS10}
{\sc Micha{\l} Adamaszek, Artur Czumaj, and Christian Sohler}: Testing monotone
  continuous distributions on high-dimensional real cubes.
\newblock In {\em Proc. 21st Ann. ACM-SIAM Symp. on Discrete Algorithms
  (SODA'10)}, pp. 56--65. ACM Press, 2010.
\newblock Short version in
  \href{http://dx.doi.org/10.1007/978-3-642-16367-8_13}{Property Testing}.
\newblock [\epfmt{acm}{1873601.1873607}]

\bibitem{AAKMRX}\bibhead{AAKMRX}
{\sc Noga Alon, Alexandr Andoni, Tali Kaufman, Kevin Matulef, Ronitt Rubinfeld,
  and Ning Xie}: Testing $k$-wise and almost $k$-wise independence.
\newblock In {\em Proc. 39th STOC}, pp. 496--505. ACM Press, 2007.
\newblock [\epfmtdoi{10.1145/1250790.1250863}]

\bibitem{ADPR}\bibhead{ADPR}
{\sc Noga Alon, Seannie Dar, Michal Parnas, and Dana Ron}: Testing of
  clustering.
\newblock {\em SIAM J. Discrete Math.}, 16(3):393--417, 2003.
\newblock Preliminary version in
  \href{http://dx.doi.org/10.1109/SFCS.2000.892111}{FOCS'00}.
\newblock [\epfmtdoi{10.1137/S0895480102410973}]

\bibitem{AlonMS99}\bibhead{AlonMS99}
{\sc Noga Alon, Yossi Matias, and Mario Szegedy}: The space complexity of
  approximating the frequency moments.
\newblock {\em J. Comput. System Sci.}, 58(1):137--147, 1999.
\newblock Preliminary version in
  \href{http://dx.doi.org/10.1145/237814.237823}{STOC'96}.
\newblock [\epfmtdoi{10.1006/jcss.1997.1545}]

\bibitem{AIOR}\bibhead{AIOR}
{\sc Alexandr Andoni, Piotr Indyk, Krzysztof Onak, and Ronitt Rubinfeld}:
  External sampling.
\newblock In {\em Proc. 36th Internat. Colloq. on Automata, Languages and
  Programming (ICALP'09)}, pp. 83--94. Springer, 2009.
\newblock [\epfmtdoi{10.1007/978-3-642-02927-1\_9}]

\bibitem{DNNR09}\bibhead{DNNR09}
{\sc Khanh~Do Ba, Huy~L. Nguyen, Huy~N. Nguyen, and Ronitt Rubinfeld}:
  Sublinear time algorithms for {E}arth {M}over's {D}istance.
\newblock {\em Theory Comput. Syst.}, 48(2):428--442, 2011.
\newblock Preprint available at \href{http://arxiv.org/abs/0904.0292}{arXiv}.
\newblock [\epfmtdoi{10.1007/s00224-010-9265-8}]

\bibitem{BJKST02}\bibhead{BJKST02}
{\sc Ziv Bar-Yossef, T.~S. Jayram, Ravi Kumar, D.~Sivakumar, and Luca
  Trevisan}: Counting distinct elements in a data stream.
\newblock In {\em Proc. 6th Internat. Workshop on Randomization and Computation
  (RANDOM'02)}, pp. 1--10. Springer, 2002.
\newblock [\epfmtdoi{10.1007/3-540-45726-7\_1}]

\bibitem{Batu}\bibhead{Batu}
{\sc Tu\v{g}kan Batu}: {\em Testing properties of distributions}.
\newblock Ph.\,D.\ thesis, Computer Science department, Cornell University,
  2001.
\newblock \url{http://www.maths.lse.ac.uk/Personal/batu/papers/t.ps}.
\newblock [\epfmt{acm}{934334}]

\bibitem{BDKR}\bibhead{BDKR}
{\sc Tu\v{g}kan Batu, Sanjoy Dasgupta, Ravi Kumar, and Ronitt Rubinfeld}: The
  complexity of approximating the entropy.
\newblock {\em SIAM J. Comput.}, 35(1):132--150, 2005.
\newblock Preliminary versions in
  \href{http://dx.doi.org/10.1109/CCC.2002.1004329}{CCC'02} and
  \href{http://dx.doi.org/10.1145/509907.510005}{STOC'02}.
\newblock [\epfmtdoi{10.1137/S0097539702403645}]

\bibitem{BFFKRW}\bibhead{BFFKRW}
{\sc Tu\v{g}kan Batu, Eldar Fischer, Lance Fortnow, Ravi Kumar, Ronitt
  Rubinfeld, and Patrick White}: Testing random variables for independence and
  identity.
\newblock In {\em Proc. 42nd FOCS}, pp. 442--451. IEEE Comp. Soc. Press, 2001.
\newblock [\epfmtdoi{10.1109/SFCS.2001.959920}]

\bibitem{BFRSW}\bibhead{BFRSW}
{\sc Tu\v{g}kan Batu, Lance Fortnow, Ronitt Rubinfeld, Warren~D. Smith, and
  Patrick White}: Testing that distributions are close.
\newblock In {\em Proc. 41st FOCS}, pp. 259--269. IEEE Comp. Soc. Press, 2000.
\newblock [\epfmtdoi{10.1109/SFCS.2000.892113}]

\bibitem{BFRSW-long}\bibhead{BFRSW-long}
{\sc Tu\v{g}kan Batu, Lance Fortnow, Ronitt Rubinfeld, Warren~D. Smith, and
  Patrick White}: Testing closeness of discrete distributions.
\newblock Technical report, 2010.
\newblock Preliminary version in
  \href{http://dx.doi.org/10.1109/SFCS.2000.892113}{FOCS'00}. Accepted for
  publication in JACM.
\newblock [\epfmt{arxiv}{1009.5397}]

\bibitem{BKR}\bibhead{BKR}
{\sc Tu\v{g}kan Batu, Ravi Kumar, and Ronitt Rubinfeld}: Sublinear algorithms
  for testing monotone and unimodal distributions.
\newblock In {\em Proc. 36th STOC}, pp. 381--390. ACM Press, 2004.
\newblock [\epfmtdoi{10.1145/1007352.1007414}]

\bibitem{BravermanOstrovsky10c}\bibhead{BravermanOstrovsky10c}
{\sc Vladimir Braverman and Rafail Ostrovsky}: Measuring $k$-wise independence
  of streaming data.
\newblock Technical report, 2008.
\newblock [\epfmt{arxiv}{0806.4790v1}]

\bibitem{BravermanOstrovsky10a}\bibhead{BravermanOstrovsky10a}
{\sc Vladimir Braverman and Rafail Ostrovsky}: Measuring independence of
  datasets.
\newblock In {\em Proc. 42nd STOC}, pp. 271--280. ACM Press, 2010.
\newblock [\epfmtdoi{10.1145/1806689.1806728}]

\bibitem{BravermanOstrovsky10b}\bibhead{BravermanOstrovsky10b}
{\sc Vladimir Braverman and Rafail Ostrovsky}: Zero-one frequency laws.
\newblock In {\em Proc. 42nd STOC}, pp. 281--290. ACM Press, 2010.
\newblock [\epfmtdoi{10.1145/1806689.1806729}]

\bibitem{CoppKum}\bibhead{CoppKum}
{\sc Don Coppersmith and Ravi Kumar}: An improved data stream algorithm for
  frequency moments.
\newblock In {\em Proc. 15th Ann. ACM-SIAM Symp. on Discrete Algorithms
  (SODA'04)}, pp. 151--156. ACM Press, 2004.
\newblock [\epfmt{acm}{982815}]

\bibitem{CDIM}\bibhead{CDIM}
{\sc Graham Cormode, Mayur Datar, Piotr Indyk, and S.~Muthukrishnan}: Comparing
  data streams using {H}amming norms (how to zero in).
\newblock {\em IEEE Trans. Knowl. Data Eng.}, 15(3):529--540, 2003.
\newblock Preliminary version in
  \href{http://www.vldb.org/conf/2002/S10P02.pdf}{VLDB'02}.
\newblock [\epfmtdoi{10.1109/TKDE.2003.1198388}]

\bibitem{Csiszar67}\bibhead{Csiszar67}
{\sc Imre Csisz\'{a}r}: Information-type measures of difference of probability
  distributions and indirect observations.
\newblock {\em Studia Scientiarum Mathematicarum Hungarica}, 2:299--318, 1967.

\bibitem{CS07}\bibhead{CS07}
{\sc Artur Czumaj and Christian Sohler}: Testing expansion in bounded-degree
  graphs.
\newblock {\em Combinatorics, Probability {\&} Computing}, 19(5-6):693--709,
  2010.
\newblock Preliminary version in
  \href{http://dx.doi.org/10.1109/FOCS.2007.33}{FOCS'07}.
\newblock [\epfmtdoi{10.1017/S096354831000012X}]

\bibitem{FeigenbaumKSV99}\bibhead{FeigenbaumKSV99}
{\sc Joan Feigenbaum, Sampath Kannan, Martin~J. Strauss, and Mahesh
  Viswanathan}: An approximate ${L}^1$-difference algorithm for massive data
  streams.
\newblock {\em SIAM J. Comput.}, 32(1):131--151, January 2003.
\newblock Preliminary verison in
  \href{http://dx.doi.org/10.1109/SFFCS.1999.814623}{FOCS'99}.
\newblock [\epfmtdoi{10.1137/S0097539799361701}]

\bibitem{feller}\bibhead{feller}
{\sc William Feller}: {\em An introduction to probability theory and its
  applications}.
\newblock Volume~1.
\newblock Wiley, 3rd edition, 1967.

\bibitem{FM}\bibhead{FM}
{\sc Eldar Fischer and Arie Matsliah}: Testing graph isomorphism.
\newblock {\em SIAM J. Comput.}, 38(1):207--225, 2008.
\newblock Preliminary version in
  \href{http://dl.acm.org/citation.cfm?id=1109591}{SODA'06}.
\newblock [\epfmtdoi{10.1137/070680795}]

\bibitem{FongS00}\bibhead{FongS00}
{\sc Jessica~H. Fong and Martin Strauss}: An approximate ${L}^p$ difference
  algorithm for massive data streams.
\newblock {\em Discrete Mathematics {\&} Theoretical Computer Science},
  4(2):301--322, 2001.
\newblock
  \href{http://www.dmtcs.org/dmtcs-ojs/index.php/dmtcs/issue/view/62}{DMTCS}.
  Preliminary version in \href{http://dx.doi.org/
  10.1007/3-540-46541-3_16}{STACS'00}.

\bibitem{GRdist}\bibhead{GRdist}
{\sc Oded Goldreich and Dana Ron}: On testing expansion in bounded-degree
  graphs.
\newblock In {\em Studies in Complexity and Cryptography}, pp. 68--75.
  Springer, 2011.
\newblock Preliminary version in
  \href{http://eccc.hpi-web.de/report/2000/020/}{ECCC}.
\newblock [\epfmtdoi{10.1007/978-3-642-22670-0\_9}]

\bibitem{GMV}\bibhead{GMV}
{\sc Sudipto Guha, Andrew McGregor, and Suresh Venkatasubramanian}: Sublinear
  estimation of entropy and information distances.
\newblock {\em ACM Trans. Algorithms}, 5(4):35, 2009.
\newblock Preliminary version in
  \href{http://dx.doi.org/10.1145/1109557.1109637}{SODA'06}.
\newblock [\epfmtdoi{10.1145/1597036.1597038}]

\bibitem{Harris75}\bibhead{Harris75}
{\sc Bernard Harris}: The statistical estimation of entropy in the
  non-parametric case.
\newblock {\em Colloquia Mathematica Societatis J\'anos Bolyai}, 16:323--355,
  1975.
\newblock
  \href{http://oai.dtic.mil/oai/oai?verb=getRecord&metadataPrefix=html&identifier=ADA020217}{DTIC
  Online}.

\bibitem{IM08}\bibhead{IM08}
{\sc Piotr Indyk and Andrew McGregor}: Declaring independence via the sketching
  of sketches.
\newblock In {\em Proc. 19th Ann. ACM-SIAM Symp. on Discrete Algorithms
  (SODA'08)}, pp. 737--745. ACM Press, 2008.
\newblock [\epfmt{acm}{1347163}]

\bibitem{IKOS09}\bibhead{IKOS09}
{\sc Yuval Ishai, Eyal Kushilevitz, Rafail Ostrovsky, and Amit Sahai}:
  Extracting correlations.
\newblock In {\em Proc. 50th FOCS}, pp. 261--270. IEEE Comp. Soc. Press, 2009.
\newblock [\epfmtdoi{10.1109/FOCS.2009.56}]

\bibitem{KS08}\bibhead{KS08}
{\sc Satyen Kale and Comandur Seshadhri}: An expansion tester for bounded
  degree graphs.
\newblock {\em SIAM J. Comput.}, 40(3):709--720, 2011.
\newblock Preliminary versions in
  \href{http://dx.doi.org/10.1007/978-3-540-70575-8_43}{ICALP'08} and
  \href{http://eccc.hpi-web.de/report/2007/076/}{ECCC}.
\newblock [\epfmtdoi{10.1137/100802980}]

\bibitem{Knuthv2}\bibhead{Knuthv2}
{\sc Donald Knuth}: {\em The Art of Computer Programming: Seminumerical
  Algorithms}.
\newblock Volume~2.
\newblock Addison Wesley, Phillipines, 1969.

\bibitem{LRR11}\bibhead{LRR11}
{\sc Reut Levi, Dana Ron, and Ronitt Rubinfeld}: Testing properties of
  collections of distributions.
\newblock In {\em Proc. 2nd Symp. Innovations in Comput. Sci. (ICS'11)}, pp.
  179--194, 2011.
\newblock
  \href{http://conference.itcs.tsinghua.edu.cn/ICS2011/content/papers/13.html}{ICS'11}.

\bibitem{Ma81}\bibhead{Ma81}
{\sc Shang-{keng} Ma}: Calculation of entropy from data of motion.
\newblock {\em J. Statistical Physics}, 26(2):221--240, 1981.
\newblock [\epfmtdoi{10.1007/BF01013169}]

\bibitem{bern}\bibhead{bern}
{\sc Dragoslav~S. Mitrinovi\'{c} and Petar~M. Vasi\'{c}}: {\em Analytic
  Inequalities}.
\newblock Springer, 1970.

\bibitem{NS08}\bibhead{NS08}
{\sc Asaf Nachmias and Asaf Shapira}: Testing the expansion of a graph.
\newblock {\em Inform. and Comput.}, 208(4):309--314, 2010.
\newblock Preliminary version in
  \href{http://eccc.hpi-web.de/report/2007/118/}{ECCC}.
\newblock [\epfmtdoi{10.1016/j.ic.2009.09.002}]

\bibitem{NemBialekDeRuyt}\bibhead{NemBialekDeRuyt}
{\sc Ilya Nemenman, William Bialek, and {Rob de Ruyter van} Steveninck}:
  Entropy and information in neural spike trains: Progress on the sampling
  problem.
\newblock {\em Phys. Rev. E}, 69(5):056111, 2004.
\newblock Preprint available at
  \href{http://arxiv.org/abs/physics/0306063}{arXiv}.
\newblock [\epfmtdoi{10.1103/PhysRevE.69.056111}]

\bibitem{Paninskient}\bibhead{Paninskient}
{\sc Liam Paninski}: Estimation of entropy and mutual information.
\newblock {\em Neural Computation}, 15(6):1191--1253, 2003.
\newblock [\epfmtdoi{10.1162/089976603321780272}]

\bibitem{Paninskisub}\bibhead{Paninskisub}
{\sc Liam Paninski}: Estimating entropy on $m$ bins given fewer than $m$
  samples.
\newblock {\em IEEE Trans. Inform. Theory}, 50(9):2200--2203, 2004.
\newblock [\epfmtdoi{10.1109/TIT.2004.833360}]

\bibitem{Paninski2}\bibhead{Paninski2}
{\sc Liam Paninski}: A coincidence-based test for uniformity given very
  sparsely sampled discrete data.
\newblock {\em IEEE Trans. Inform. Theory}, 54(10):4750--4755, 2008.
\newblock [\epfmtdoi{10.1109/TIT.2008.928987}]

\bibitem{RRRS07}\bibhead{RRRS07}
{\sc Sofya Raskhodnikova, Dana Ron, Ronitt Rubinfeld, and Adam Smith}:
  Sublinear algorithms for approximating string compressibility.
\newblock {\em Algorithmica}, 65(3):685--709, 2013.
\newblock Preliminary version in
  \href{http://dx.doi.org/10.1007/978-3-540-74208-1_44}{RANDOM'07}.
\newblock [\epfmtdoi{10.1007/s00453-012-9618-6}]

\bibitem{RRSS}\bibhead{RRSS}
{\sc Sofya Raskhodnikova, Dana Ron, Amir Shpilka, and Adam Smith}: Strong lower
  bonds for approximating distributions support size and the distinct elements
  problem.
\newblock {\em SIAM J. Comput.}, 39(3):813--842, 2009.
\newblock Preliminary version in
  \href{http://dx.doi.org/10.1109/FOCS.2007.47}{FOCS'07}.
\newblock [\epfmtdoi{10.1137/070701649}]

\bibitem{roos}\bibhead{roos}
{\sc Bero Roos}: On the rate of multivariate {P}oisson convergence.
\newblock {\em J. Multivar. Anal.}, 69(1):120--134, 1999.
\newblock [\epfmtdoi{10.1006/jmva.1998.1789}]

\bibitem{RServ}\bibhead{RServ}
{\sc Ronitt Rubinfeld and Rocco~A. Servedio}: Testing monotone high-dimensional
  distributions.
\newblock {\em Random Structures {\&} Algorithms}, 34(1):24--44, 2009.
\newblock Preliminary version in
  \href{http://dx.doi.org/10.1145/1060590.1060613}{STOC'05}.
\newblock [\epfmtdoi{10.1002/rsa.20247}]

\bibitem{RX10}\bibhead{RX10}
{\sc Ronitt Rubinfeld and Ning Xie}: Robust characterizations of $k$-wise
  independence over product spaces and related testing results.
\newblock {\em Random Structures \& Algorithms}, 2012 (online).
\newblock Preliminary version in
  \href{http://dx.doi.org/10.1007/978-3-642-14165-2\_48}{ICALP 2010}.
\newblock [\epfmtdoi{10.1002/rsa.20423}]

\bibitem{Bialek98}\bibhead{Bialek98}
{\sc Steve~P. Strong, Roland Koberle, {Rob R. de Ruyter van} Steveninck, and
  William Bialek}: Entropy and information in neural spike trains.
\newblock {\em Phys. Rev. Lett.}, 80(1):197--200, 1998.
\newblock [\epfmtdoi{10.1103/PhysRevLett.80.197}]

\bibitem{szpankowski}\bibhead{szpankowski}
{\sc Wojciech Szpankowski}: {\em Average Case Analysis of Algorithms on
  Sequences}.
\newblock John Wiley {\&} Sons, New York, 2001.
\newblock [\epfmt{acm}{517168}]

\bibitem{valphd}\bibhead{valphd}
{\sc Paul Valiant}: {\em Testing symmetric properties of distributions}.
\newblock Ph.\,D.\ thesis, CSAIL, MIT, 2008.
\newblock \href{http://hdl.handle.net/1721.1/44717}{DSpace@MIT}.

\bibitem{val}\bibhead{val}
{\sc Paul Valiant}: Testing symmetric properties of distributions.
\newblock {\em SIAM J. Comput.}, 40(6):1927--1968, 2011.
\newblock Preliminary version in
  \href{http://dx.doi.org/10.1145/1374376.1374432}{STOC'08}.
\newblock [\epfmtdoi{10.1137/080734066}]

\bibitem{Patrick-phd}\bibhead{Patrick-phd}
{\sc Patrick White}: Testing random variables for independence and identity.
\newblock Unpublished manuscript, 2002.

\bibitem{WolpertW95}\bibhead{WolpertW95}
{\sc David~H. Wolpert and David~R. Wolf}: Estimating functions of probability
  distributions from a finite set of samples.
\newblock {\em Physical Review E}, 52(6):6841--6854, 1995.
\newblock [\epfmtdoi{10.1103/PhysRevE.52.6841}]

\bibitem{Yamanishi95}\bibhead{Yamanishi95}
{\sc Kenji Yamanishi}: Probably almost discriminative learning.
\newblock {\em Machine Learning}, 18(1):23--50, 1995.
\newblock Preliminary version in
  \href{http://dx.doi.org/10.1145/130385.130404}{COLT'92}.
\newblock [\epfmtdoi{10.1007/BF00993820}]

\end{thebibliography}

